
\chapwithtoc{Závěr}

V této práci jsme se zabývali zpřístupněním generátorů \texttt{geng} a
\texttt{plantri} z moderního C++ kódu. Navrhli a implementovali jsme knihovnu
PlanG, která nad těmito nástroji poskytuje jednotné rozhraní pro konfiguraci
generování, filtrování a zpracování grafů.

Součástí řešení je \texttt{GraphView}, které zpřístupňuje aktuální graf bez
automatického kopírování do nové datové struktury. Pro oba backendy jsme
připravili rozhraní použitelné s vybranými koncepty Boost Graph Library.
Generátor \texttt{geng} jsme navíc rozšířili o podporu barevných tříd a
zakotvených vrcholů.

Výsledné řešení jsme ověřili funkčními testy a porovnali s původními
generátory. Měření ukázala, že při čistém generování dosahuje PlanG výkonu
srovnatelného s nativními programy. PlanG tak umožňuje používat výkonné
generátory grafů pohodlněji z C++ programu, aniž by bylo nutné každý graf
automaticky kopírovat nebo psát samostatné pluginy.

\section*{Možná rozšíření}

Dalším směrem rozvoje je podpora dalších generátorů grafů a rozšíření rozhraní
o další globální funkce a algoritmy, které pracují přímo nad objekty
\texttt{GraphView}.